\begin{bt}%[9D6H18-2]
% nguon: Tuyển tập đề tham khảo môn Toán Tuyển sinh 10 năm học 2025-2026 TP.HCM – Đề tham khảo Huyện Bình Chánh (Đề 1) – Bài 1
Cho hàm số $y=\dfrac{x^2}{4}$ có đồ thị là Parabol $(P)$.
\begin{enumerate}[a)]
	\item Vẽ đồ thị $(P)$ của hàm số trên hệ trục toạ độ $Oxy$.
	\item Tìm toạ độ những điểm $M$ thuộc $(P)$ có tung độ bằng $2$ lần hoành độ.
\end{enumerate}
\loigiai{
\begin{enumerate}[a)]
	\item Bảng giá trị tương ứng của $x$ và $y$:
	\begin{center}
	\begin{tabular}{|c|c|c|c|c|c|}
	\hline
	$x$ & $-4$ & $-2$ & $0$ & $2$ & $4$ \\ \hline
	$y=\dfrac{x^2}{4}$ & $4$ & $1$ & $0$ & $1$ & $4$ \\ \hline
	\end{tabular}
	\end{center}
	Đồ thị $(P)$ là một đường parabol đi qua gốc toạ độ $O(0;0)$, nhận trục tung $Oy$ làm trục đối xứng và quay bề lõm lên phía trên.
	\item Gọi $M(x_0; y_0)$ là điểm thuộc $(P)$.\\
	Theo giả thiết, tung độ bằng $2$ lần hoành độ nên $y_0 = 2x_0$.\\
	Vì $M \in (P)$ nên:
	$$ \dfrac{x_0^2}{4} = 2x_0 \iff x_0^2 - 8x_0 = 0 \iff x_0(x_0 - 8) = 0 \iff \left[\begin{aligned} x_0 &= 0 \\ x_0 &= 8. \end{aligned}\right. $$
	Với $x_0 = 0 \implies y_0 = 0$, ta được điểm $O(0;0)$.\\
	Với $x_0 = 8 \implies y_0 = 2 \cdot 8 = 16$, ta được điểm $M(8;16)$.\\
	Vậy có hai điểm thoả mãn là $O(0;0)$ và $M(8;16)$.
\end{enumerate}
}
\end{bt}

\begin{bt}%[9D6H20-2]
% nguon: Tuyển tập đề tham khảo môn Toán Tuyển sinh 10 năm học 2025-2026 TP.HCM – Đề tham khảo Huyện Bình Chánh (Đề 1) – Bài 2
Cho phương trình bậc hai: $3x^2-12x+2=0$.
\begin{enumerate}[a)]
	\item Chứng minh phương trình trên luôn có hai nghiệm phân biệt $x_1, x_2$.
	\item Không giải phương trình, hãy tính giá trị của biểu thức:
	$$ A = x_1^3 + x_2^3 - (x_1+x_2). $$
\end{enumerate}
\loigiai{
\begin{enumerate}[a)]
	\item Phương trình $3x^2-12x+2=0$ có các hệ số $a=3$, $b'=-6$, $c=2$.\\
	Biệt thức thu gọn:
	$$ \Delta' = (-6)^2 - 3 \cdot 2 = 36 - 6 = 30 > 0. $$
	Vì $\Delta' > 0$ nên phương trình luôn có hai nghiệm phân biệt $x_1, x_2$.
	\item Theo hệ thức Vi-ét, ta có:
	$$ \begin{cases} S = x_1+x_2 = -\dfrac{-12}{3} = 4 \\ P = x_1 x_2 = \dfrac{2}{3}. \end{cases} $$
	Biến đổi tổng các lập phương theo $S$ và $P$:
	$$ x_1^3 + x_2^3 = (x_1+x_2)^3 - 3x_1 x_2(x_1+x_2) = S^3 - 3PS. $$
	Do đó biểu thức $A$ trở thành:
	$$ A = S^3 - 3PS - S. $$
	Thay $S = 4$ và $P = \dfrac{2}{3}$ vào biểu thức $A$, ta được:
	$$ A = 4^3 - 3 \cdot \dfrac{2}{3} \cdot 4 - 4 = 64 - 8 - 4 = 52. $$
\end{enumerate}
}
\end{bt}

\begin{bt}%[9D6V21-3]
% nguon: Tuyển tập đề tham khảo môn Toán Tuyển sinh 10 năm học 2025-2026 TP.HCM – Đề tham khảo Huyện Bình Chánh (Đề 1) – Bài 3
Một trường học xây dựng một sân bóng rổ hình chữ nhật có chiều dài $26\text{ m}$ và chiều rộng $14\text{ m}$. Theo thiết kế, người ta xây dựng thêm một lối đi có dạng chữ L bao quanh hai cạnh vuông góc của sân bóng rổ. Gọi $x\text{ (m)}$ là bề rộng của lối đi ($x > 0$).
\begin{enumerate}[a)]
	\item Viết biểu thức $S$ biểu diễn diện tích của lối đi theo $x$.
	\item Bạn An đi bộ từ cửa vào đến cửa ra dọc theo viền ngoài của lối đi (tổng độ dài hai cạnh ngoài của lối đi). Tính quãng đường bạn An đã đi biết diện tích của lối đi theo thiết kế là $129\text{ m}^2$.
\end{enumerate}
\loigiai{
\begin{enumerate}[a)]
	\item Lối đi bao gồm hai dải hình chữ nhật dọc theo chiều dài $26\text{ m}$ và chiều rộng $14\text{ m}$ cùng một góc hình vuông kích thước $x \times x$ ở chỗ giao nhau.\\
	Do đó biểu thức diện tích của lối đi theo $x$ là:
	$$ S(x) = 26x + 14x + x^2 = x^2 + 40x\text{ (m}^2\text{)}. $$
	\item Theo thiết kế diện tích lối đi là $129\text{ m}^2$, ta có phương trình:
	$$ x^2 + 40x = 129 \iff x^2 + 40x - 129 = 0. $$
	Giải phương trình bậc hai:
	$$ \Delta' = 20^2 - 1 \cdot (-129) = 400 + 129 = 529 = 23^2. $$
	Phương trình có hai nghiệm:
	$$ x_1 = -20 + 23 = 3\text{ (nhận do } x > 0\text{)}, \quad x_2 = -20 - 23 = -43\text{ (loại)}. $$
	Bề rộng của lối đi là $x = 3\text{ m}$.\\
	Hai cạnh viền ngoài của lối đi có kích thước lần lượt là $26 + x = 26 + 3 = 29\text{ m}$ và $14 + x = 14 + 3 = 17\text{ m}$.\\
	Quãng đường bạn An đã đi là:
	$$ 29 + 17 = 46\text{ (m)}. $$
\end{enumerate}
}
\end{bt}

\begin{bt}%[9H0V31-3]
% nguon: Tuyển tập đề tham khảo môn Toán Tuyển sinh 10 năm học 2025-2026 TP.HCM – Đề tham khảo Huyện Bình Chánh (Đề 1) – Bài 4
Nước giải khát thường đựng trong lon nhôm hình trụ. Cỡ lon phổ biến có chiều cao $12\text{ cm}$, đường kính đáy $6{,}5\text{ cm}$ và chứa được khoảng $335\text{ ml}$ chất lỏng. Hiện nay các nhà sản xuất có xu hướng thiết kế lon nhôm kiểu dáng thon cao hơn với chiều cao $14\text{ cm}$ và đường kính đáy $6\text{ cm}$.
\begin{enumerate}[a)]
	\item Một lon nước ngọt kiểu mới cao $14\text{ cm}$, đường kính đáy $6\text{ cm}$ có thể chứa hết lượng nước ngọt của một lon cỡ phổ biến ($335\text{ ml}$) không? Vì sao? (lấy $\pi \approx 3{,}14$).
	\item Biết chi phí sản xuất lon tỉ lệ thuận với diện tích toàn phần của lon hình trụ ($S_{tp} = \pi d h + 2\pi r^2$). Hỏi chi phí sản xuất vỏ lon kiểu dáng mới tăng khoảng bao nhiêu phần trăm so với lon cỡ phổ biến? (kết quả làm tròn đến hàng phần mười).
\end{enumerate}
\loigiai{
\begin{enumerate}[a)]
	\item Bán kính đáy của lon nước ngọt kiểu mới là:
	$$ r_1 = \dfrac{6}{2} = 3\text{ cm}. $$
	Thể tích của lon kiểu mới là:
	$$ V_1 = \pi r_1^2 h_1 = 3{,}14 \cdot 3^2 \cdot 14 = 395{,}64\text{ cm}^3 = 395{,}64\text{ ml}. $$
	Vì $395{,}64\text{ ml} > 335\text{ ml}$ nên lon nước ngọt kiểu mới hoàn toàn có thể chứa hết lượng nước ngọt của lon cỡ phổ biến.
	\item Tính diện tích toàn phần của hai loại lon:
	\begin{itemize}
		\item Đối với lon cỡ phổ biến ($h_0 = 12\text{ cm}, r_0 = \dfrac{6{,}5}{2} = 3{,}25\text{ cm}$):
		$$ S_{tp0} = 2\pi r_0 h_0 + 2\pi r_0^2 = 2\pi \cdot 3{,}25 \cdot 12 + 2\pi \cdot (3{,}25)^2 = 78\pi + 21{,}125\pi = 99{,}125\pi\text{ (cm}^2\text{)}. $$
		\item Đối với lon kiểu mới ($h_1 = 14\text{ cm}, r_1 = 3\text{ cm}$):
		$$ S_{tp1} = 2\pi r_1 h_1 + 2\pi r_1^2 = 2\pi \cdot 3 \cdot 14 + 2\pi \cdot 3^2 = 84\pi + 18\pi = 102\pi\text{ (cm}^2\text{)}. $$
	\end{itemize}
	Tỉ lệ tăng chi phí sản xuất vỏ lon tương ứng với tỉ lệ tăng diện tích toàn phần:
	$$ \dfrac{S_{tp1} - S_{tp0}}{S_{tp0}} \cdot 100\% = \dfrac{102\pi - 99{,}125\pi}{99{,}125\pi} \cdot 100\% = \dfrac{2{,}875}{99{,}125} \cdot 100\% \approx 2{,}9\%. $$
	Vậy chi phí sản xuất lon nước ngọt kiểu mới tăng khoảng $2{,}9\%$ so với lon cỡ phổ biến.
\end{enumerate}
}
\end{bt}

\begin{bt}%[9D1H3-4]
% nguon: Tuyển tập đề tham khảo môn Toán Tuyển sinh 10 năm học 2025-2026 TP.HCM – Đề tham khảo Huyện Bình Chánh (Đề 1) – Bài 5
Có hai loại can nhựa dùng để đựng hoá chất. Nếu lấy lượng hoá chất chứa đầy trong $2$ can loại lớn đổ sang can loại bé thì rót đầy được $4$ can bé và còn dư ra $2$ lít. Nếu lấy lượng hoá chất chứa đầy trong $7$ can loại bé đổ sang can loại lớn thì rót đầy được $3$ can lớn và còn dư ra $1$ lít. Tìm dung tích của mỗi loại can nhựa đó.
\loigiai{
Gọi dung tích của một can loại lớn và một can loại bé lần lượt là $x$ và $y$ ($x > y > 0$, đơn vị: lít).\\
Theo trường hợp thứ nhất: lượng hoá chất trong $2$ can lớn đổ đầy được $4$ can bé và dư $2$ lít, ta có:
$$ 2x = 4y + 2 \iff x - 2y = 1. $$
Theo trường hợp thứ hai: lượng hoá chất trong $7$ can bé đổ đầy được $3$ can lớn và dư $1$ lít, ta có:
$$ 7y = 3x + 1 \iff 3x - 7y = -1. $$
Ta có hệ phương trình bậc nhất hai ẩn:
$$ \begin{cases} x - 2y = 1 \\ 3x - 7y = -1 \end{cases} \iff \begin{cases} x = 2y + 1 \\ 3(2y + 1) - 7y = -1 \end{cases} \iff \begin{cases} -y = -4 \\ x = 2y + 1 \end{cases} \iff \begin{cases} y = 4\text{ (thoả mãn)} \\ x = 2 \cdot 4 + 1 = 9\text{ (thoả mãn)}. \end{cases} $$
Vậy dung tích của can loại lớn là $9$ lít và dung tích của can loại bé là $4$ lít.
}
\end{bt}

\begin{bt}%[9D8H26-3]
% nguon: Tuyển tập đề tham khảo môn Toán Tuyển sinh 10 năm học 2025-2026 TP.HCM – Đề tham khảo Huyện Bình Chánh (Đề 1) – Bài 6
Nhóm học sinh tình nguyện khối 9 của một trường THCS gồm có $6$ bạn: $3$ bạn nam là Trung (lớp 9A), Quý (lớp 9A), Việt (lớp 9C); và $3$ bạn nữ là An (lớp 9A), Châu (lớp 9B), Hương (lớp 9D). Nhà trường chọn ngẫu nhiên $1$ bạn trong nhóm để đại diện tham gia hoạt động tình nguyện.
\begin{enumerate}[a)]
	\item Liệt kê tất cả các kết quả có thể xảy ra của phép thử trên. Phép thử có tất cả bao nhiêu kết quả có thể xảy ra?
	\item Tính xác suất của các biến cố sau:
	\begin{itemize}
		\item Biến cố $A$: ``Bạn được chọn ra là học sinh nữ''.
		\item Biến cố $B$: ``Bạn được chọn ra là học sinh thuộc lớp 9A''.
	\end{itemize}
\end{enumerate}
\loigiai{
\begin{enumerate}[a)]
	\item Tập hợp tất cả các kết quả có thể xảy ra của phép thử là:
	$$ \Omega = \{\text{Trung}, \text{Quý}, \text{Việt}, \text{An}, \text{Châu}, \text{Hương}\}. $$
	Phép thử có tất cả $6$ kết quả có thể xảy ra, tức là $n(\Omega) = 6$.
	\item Tính xác suất của từng biến cố:
	\begin{itemize}
		\item Biến cố $A$ có $3$ kết quả thuận lợi là $\{\text{An}, \text{Châu}, \text{Hương}\}$, suy ra $n(A) = 3$.\\
		Xác suất của biến cố $A$ là:
		$$ P(A) = \dfrac{n(A)}{n(\Omega)} = \dfrac{3}{6} = \dfrac{1}{2} = 0{,}5 \ (50\%). $$
		\item Các bạn học sinh thuộc lớp 9A gồm Trung, Quý, An nên biến cố $B$ có $3$ kết quả thuận lợi, suy ra $n(B) = 3$.\\
		Xác suất của biến cố $B$ là:
		$$ P(B) = \dfrac{n(B)}{n(\Omega)} = \dfrac{3}{6} = \dfrac{1}{2} = 0{,}5 \ (50\%). $$
	\end{itemize}
\end{enumerate}
}
\end{bt}

\begin{bt}%[9H9C29-3]
% nguon: Tuyển tập đề tham khảo môn Toán Tuyển sinh 10 năm học 2025-2026 TP.HCM – Đề tham khảo Huyện Bình Chánh (Đề 1) – Bài 7
Cho điểm $M$ nằm ngoài đường tròn $(O;R)$ sao cho $OM = 2R$. Vẽ hai tiếp tuyến $MA, MB$ với đường tròn $(O)$ ($A, B$ là hai tiếp điểm). Gọi $H$ là giao điểm của $MO$ và đoạn thẳng $AB$. Vẽ đường kính $AC$ của đường tròn $(O)$.
\begin{enumerate}[a)]
	\item Chứng minh đường thẳng $MO$ vuông góc với $AB$ tại $H$ và bốn điểm $M, A, O, B$ cùng thuộc một đường tròn.
	\item Chứng minh $BC \parallel MO$ và $OH \cdot OM = \dfrac{AC^2}{4}$.
	\item Tính độ dài đoạn thẳng $AB$ theo $R$.
\end{enumerate}
\loigiai{
\begin{enumerate}[a)]
	\item Theo tính chất của hai tiếp tuyến cắt nhau, ta có $MA = MB$ và $OA = OB = R$, nên $MO$ là đường trung trực của đoạn thẳng $AB \implies MO \perp AB$ tại trung điểm $H$ của $AB$.\\
	Vì $MA, MB$ là hai tiếp tuyến nên $\widehat{MAO} = 90^\circ$ và $\widehat{MBO} = 90^\circ$.\\
	Tứ giác $MAOB$ có $\widehat{MAO} + \widehat{MBO} = 90^\circ + 90^\circ = 180^\circ$, suy ra bốn điểm $M, A, O, B$ cùng thuộc đường tròn đường kính $MO$.
	\item Tam giác $ABC$ nội tiếp đường tròn $(O)$ có đường kính $AC$ nên $\widehat{ABC} = 90^\circ \implies BC \perp AB$.\\
	Lại có $MO \perp AB$ tại $H$, suy ra $BC \parallel MO$ (vì cùng vuông góc với $AB$).\\
	Xét tam giác vuông $MAO$ vuông tại $A$ có đường cao $AH \perp MO$, áp dụng hệ thức lượng ta có:
	$$ OH \cdot OM = OA^2. $$
	Vì $AC$ là đường kính nên $AC = 2R \implies OA = R = \dfrac{AC}{2} \implies OA^2 = \dfrac{AC^2}{4}$.\\
	Do đó:
	$$ OH \cdot OM = \dfrac{AC^2}{4}. $$
	\item Trong tam giác vuông $MAO$ vuông tại $A$ với $OA = R$ và $OM = 2R$, ta có:
	$$ MA = \sqrt{OM^2 - OA^2} = \sqrt{(2R)^2 - R^2} = R\sqrt{3}. $$
	Áp dụng hệ thức lượng cho đường cao $AH$ trong tam giác vuông $MAO$:
	$$ AH = \dfrac{MA \cdot OA}{OM} = \dfrac{R\sqrt{3} \cdot R}{2R} = \dfrac{R\sqrt{3}}{2}. $$
	Vì $H$ là trung điểm của dây cung $AB$ nên:
	$$ AB = 2AH = 2 \cdot \dfrac{R\sqrt{3}}{2} = R\sqrt{3}. $$
\end{enumerate}
}
\end{bt}

\begin{bt}%[9D6H18-2]
% nguon: Tuyển tập đề tham khảo môn Toán Tuyển sinh 10 năm học 2025-2026 TP.HCM – Đề tham khảo Huyện Bình Chánh (Đề 2) – Bài 1
Cho hàm số $y=-\dfrac{1}{2}x^2$ có đồ thị là Parabol $(P)$.
\begin{enumerate}[a)]
	\item Vẽ đồ thị $(P)$ trên hệ trục toạ độ $Oxy$.
	\item Tìm toạ độ những điểm trên $(P)$ sao cho hoành độ bằng $2$ lần tung độ.
\end{enumerate}
\loigiai{
\begin{enumerate}[a)]
	\item Bảng giá trị tương ứng của $x$ và $y$:
	\begin{center}
	\begin{tabular}{|c|c|c|c|c|c|}
	\hline
	$x$ & $-4$ & $-2$ & $0$ & $2$ & $4$ \\ \hline
	$y=-\dfrac{1}{2}x^2$ & $-8$ & $-2$ & $0$ & $-2$ & $-8$ \\ \hline
	\end{tabular}
	\end{center}
	Đồ thị $(P)$ là một đường parabol đi qua gốc toạ độ $O(0;0)$, nhận trục tung $Oy$ làm trục đối xứng và quay bề lõm xuống phía dưới.
	\item Gọi $A(x_0; y_0)$ là điểm thuộc $(P)$. Theo giả thiết, hoành độ bằng $2$ lần tung độ nên $x_0 = 2y_0 \iff y_0 = \dfrac{1}{2}x_0$.\\
	Vì $A \in (P)$ nên ta có:
	$$ -\dfrac{1}{2}x_0^2 = \dfrac{1}{2}x_0 \iff x_0^2 + x_0 = 0 \iff x_0(x_0 + 1) = 0 \iff \left[\begin{aligned} x_0 &= 0 \\ x_0 &= -1. \end{aligned}\right. $$
	Với $x_0 = 0 \implies y_0 = 0$, ta được điểm $O(0;0)$.\\
	Với $x_0 = -1 \implies y_0 = -\dfrac{1}{2}$, ta được điểm $A\left(-1; -\dfrac{1}{2}\right)$.\\
	Vậy có hai điểm thoả mãn là $O(0;0)$ và $A\left(-1; -\dfrac{1}{2}\right)$.
\end{enumerate}
}
\end{bt}

\begin{bt}%[9D6H20-2]
% nguon: Tuyển tập đề tham khảo môn Toán Tuyển sinh 10 năm học 2025-2026 TP.HCM – Đề tham khảo Huyện Bình Chánh (Đề 2) – Bài 2
Cho phương trình bậc hai: $3x^2+x-4=0$.
\begin{enumerate}[a)]
	\item Chứng minh phương trình trên có hai nghiệm phân biệt $x_1, x_2$.
	\item Không giải phương trình, hãy tính giá trị của biểu thức:
	$$ B = \dfrac{x_1-1}{x_2} + \dfrac{x_2-1}{x_1}. $$
\end{enumerate}
\loigiai{
\begin{enumerate}[a)]
	\item Phương trình $3x^2+x-4=0$ có các hệ số $a=3$, $b=1$, $c=-4$.\\
	Vì $a \cdot c = 3 \cdot (-4) = -12 < 0$ nên phương trình luôn có hai nghiệm phân biệt $x_1, x_2$.
	\item Theo hệ thức Vi-ét, ta có:
	$$ \begin{cases} S = x_1+x_2 = -\dfrac{1}{3} \\ P = x_1 x_2 = -\dfrac{4}{3}. \end{cases} $$
	Quy đồng mẫu thức biểu thức $B$:
	$$ B = \dfrac{x_1(x_1-1) + x_2(x_2-1)}{x_1 x_2} = \dfrac{x_1^2 + x_2^2 - (x_1+x_2)}{x_1 x_2} = \dfrac{(x_1+x_2)^2 - 2x_1 x_2 - (x_1+x_2)}{x_1 x_2}. $$
	Thay $S = -\dfrac{1}{3}$ và $P = -\dfrac{4}{3}$ vào biểu thức $B$, ta được:
	$$ B = \dfrac{\left(-\dfrac{1}{3}\right)^2 - 2\left(-\dfrac{4}{3}\right) - \left(-\dfrac{1}{3}\right)}{-\dfrac{4}{3}} = \dfrac{\dfrac{1}{9} + \dfrac{8}{3} + \dfrac{1}{3}}{-\dfrac{4}{3}} = \dfrac{\dfrac{28}{9}}{-\dfrac{4}{3}} = -\dfrac{7}{3}. $$
\end{enumerate}
}
\end{bt}

\begin{bt}%[9D1H3-4]
% nguon: Tuyển tập đề tham khảo môn Toán Tuyển sinh 10 năm học 2025-2026 TP.HCM – Đề tham khảo Huyện Bình Chánh (Đề 2) – Bài 3
Một bức tường hình thang vuông có chiều cao là $h\text{ (m)}$, đáy bé là $4x\text{ (m)}$ và đáy lớn là $8x\text{ (m)}$. Trên bức tường đó người ta thiết kế một cửa sổ hình chữ nhật có chiều rộng là $2x\text{ (m)}$ và chiều cao là $1\text{ m}$.
\begin{enumerate}[a)]
	\item Viết biểu thức biểu thị diện tích bức tường theo $x$ và $h$ (không tính phần cửa sổ).
	\item Tính diện tích của bức tường khi biết $x = 3\text{ m}$ và $h = 4\text{ m}$.
\end{enumerate}
\loigiai{
\begin{enumerate}[a)]
	\item Diện tích của toàn bộ bức tường hình thang vuông ban đầu là:
	$$ S_{\text{thang}} = \dfrac{(4x + 8x) \cdot h}{2} = \dfrac{12x \cdot h}{2} = 6xh\text{ (m}^2\text{)}. $$
	Diện tích của cửa sổ hình chữ nhật là:
	$$ S_{\text{cửa sổ}} = 2x \cdot 1 = 2x\text{ (m}^2\text{)}. $$
	Biểu thức diện tích phần tường còn lại (không tính cửa sổ) là:
	$$ S = 6xh - 2x\text{ (m}^2\text{)}. $$
	\item Thay $x = 3\text{ m}$ và $h = 4\text{ m}$ vào biểu thức diện tích, ta được:
	$$ S = 6 \cdot 3 \cdot 4 - 2 \cdot 3 = 72 - 6 = 66\text{ (m}^2\text{)}. $$
	Vậy diện tích bức tường là $66\text{ m}^2$.
\end{enumerate}
}
\end{bt}

\begin{bt}%[9H0H31-3]
% nguon: Tuyển tập đề tham khảo môn Toán Tuyển sinh 10 năm học 2025-2026 TP.HCM – Đề tham khảo Huyện Bình Chánh (Đề 2) – Bài 4
Khi thả chìm hoàn toàn một viên xúc xắc hình lập phương vào một chiếc ly đựng nước có dạng hình trụ, người ta quan sát thấy mực nước trong ly dâng lên thêm $0{,}5\text{ cm}$ và nước không bị tràn ra ngoài. Biết diện tích mặt đáy của ly nước hình trụ là $250\text{ cm}^2$. Hỏi độ dài cạnh của viên xúc xắc đó dài bao nhiêu xăng-ti-mét?
\loigiai{
Thể tích nước dâng lên trong ly hình trụ chính bằng thể tích của viên xúc xắc hình lập phương chìm trong nước.\\
Thể tích nước dâng lên là:
$$ V = S_{\text{đáy}} \cdot h_{\text{dâng}} = 250 \cdot 0{,}5 = 125\text{ (cm}^3\text{)}. $$
Gọi độ dài cạnh của viên xúc xắc hình lập phương là $a\text{ (cm)}$ ($a > 0$).\\
Ta có công thức thể tích hình lập phương là $V = a^3$, suy ra:
$$ a^3 = 125 \iff a = \sqrt[3]{125} = 5\text{ (cm)}. $$
Vậy độ dài cạnh của viên xúc xắc là $5\text{ cm}$.
}
\end{bt}

\begin{bt}%[9D1H3-4]
% nguon: Tuyển tập đề tham khảo môn Toán Tuyển sinh 10 năm học 2025-2026 TP.HCM – Đề tham khảo Huyện Bình Chánh (Đề 2) – Bài 5
Một công ty địa ốc xây dựng một toà chung cư cao cấp gồm $100$ căn hộ để bán thuộc hai loại: loại I (căn hộ $1$ phòng ngủ) có giá bán ban đầu là $1{,}7$ tỉ đồng/căn; loại II (căn hộ $2$ phòng ngủ) có giá bán ban đầu là $2{,}6$ tỉ đồng/căn. Do thay đổi chiến lược kinh doanh, chủ đầu tư quyết định điều chỉnh giá: tăng $10\%$ giá bán đối với căn hộ loại I và giảm $5\%$ giá bán đối với căn hộ loại II. Sau khi bán hết toàn bộ $100$ căn hộ, công ty thu về tổng số tiền là $211$ tỉ đồng. Hỏi toà chung cư đó có bao nhiêu căn hộ loại I và bao nhiêu căn hộ loại II?
\loigiai{
Gọi số căn hộ loại I và loại II của toà chung cư lần lượt là $x$ và $y$ ($x, y \in \mathbb{N}^*$, căn hộ).\\
Tổng số căn hộ của toà chung cư là $100$ căn, ta có phương trình:
$$ x + y = 100. $$
Giá bán thực tế của mỗi căn hộ sau khi điều chỉnh là:
\begin{itemize}
	\item Căn hộ loại I: $1{,}7 \cdot (1 + 0{,}10) = 1{,}87$ (tỉ đồng/căn).
	\item Căn hộ loại II: $2{,}6 \cdot (1 - 0{,}05) = 2{,}47$ (tỉ đồng/căn).
\end{itemize}
Tổng số tiền thu về khi bán hết toàn bộ căn hộ là $211$ tỉ đồng, ta có phương trình:
$$ 1{,}87x + 2{,}47y = 211. $$
Ta có hệ phương trình bậc nhất hai ẩn:
$$ \begin{cases} x + y = 100 \\ 1{,}87x + 2{,}47y = 211 \end{cases} \iff \begin{cases} x = 100 - y \\ 1{,}87(100 - y) + 2{,}47y = 211 \end{cases} \iff \begin{cases} 0{,}6y = 24 \\ x = 100 - y \end{cases} \iff \begin{cases} y = 40\text{ (thoả mãn)} \\ x = 60\text{ (thoả mãn)}. \end{cases} $$
Vậy toà chung cư có $60$ căn hộ loại I và $40$ căn hộ loại II.
}
\end{bt}

\begin{bt}%[9D8H26-3]
% nguon: Tuyển tập đề tham khảo môn Toán Tuyển sinh 10 năm học 2025-2026 TP.HCM – Đề tham khảo Huyện Bình Chánh (Đề 2) – Bài 6
Hộp thứ nhất chứa $1$ quả bóng màu xanh và $1$ quả bóng màu đỏ. Hộp thứ hai chứa $1$ quả bóng màu vàng và $1$ quả bóng màu đỏ. Người ta lấy ngẫu nhiên từ mỗi hộp ra $1$ quả bóng.
\begin{enumerate}[a)]
	\item Liệt kê không gian mẫu và xác định số kết quả có thể xảy ra của phép thử.
	\item Tính xác suất của các biến cố sau:
	\begin{itemize}
		\item Biến cố $M$: ``Hai quả bóng lấy ra có cùng màu''.
		\item Biến cố $N$: ``Hai quả bóng lấy ra có màu khác nhau''.
		\item Biến cố $Q$: ``Có ít nhất một quả bóng màu đỏ trong hai quả bóng lấy ra''.
	\end{itemize}
\end{enumerate}
\loigiai{
\begin{enumerate}[a)]
	\item Ký hiệu các màu bóng: X (Xanh), Đ (Đỏ), V (Vàng).\\
	Không gian mẫu của phép thử là tập hợp các cặp màu bóng lấy ra lần lượt từ hộp thứ nhất và hộp thứ hai:
	$$ \Omega = \{(\text{X}, \text{V}), (\text{X}, \text{Đ}), (\text{Đ}, \text{V}), (\text{Đ}, \text{Đ})\}. $$
	Số kết quả có thể xảy ra là $n(\Omega) = 4$.
	\item Tính xác suất của từng biến cố:
	\begin{itemize}
		\item Biến cố $M$ là hai bóng cùng màu: chỉ có kết quả $(\text{Đ}, \text{Đ})$, suy ra $n(M) = 1$.\\
		Xác suất của biến cố $M$ là:
		$$ P(M) = \dfrac{1}{4} = 0{,}25 \ (25\%). $$
		\item Biến cố $N$ là hai bóng khác màu: gồm $\{(\text{X}, \text{V}), (\text{X}, \text{Đ}), (\text{Đ}, \text{V})\}$, suy ra $n(N) = 3$.\\
		Xác suất của biến cố $N$ là:
		$$ P(N) = \dfrac{3}{4} = 0{,}75 \ (75\%). $$
		\item Biến cố $Q$ là có ít nhất một bóng màu đỏ: gồm $\{(\text{X}, \text{Đ}), (\text{Đ}, \text{V}), (\text{Đ}, \text{Đ})\}$, suy ra $n(Q) = 3$.\\
		Xác suất của biến cố $Q$ là:
		$$ P(Q) = \dfrac{3}{4} = 0{,}75 \ (75\%). $$
	\end{itemize}
\end{enumerate}
}
\end{bt}

\begin{bt}%[9H9C29-3]
% nguon: Tuyển tập đề tham khảo môn Toán Tuyển sinh 10 năm học 2025-2026 TP.HCM – Đề tham khảo Huyện Bình Chánh (Đề 2) – Bài 7
Cho đường tròn $(O; R)$ với $R = 3\text{ cm}$. Hai điểm $B, C$ thuộc $(O)$ sao cho góc ở tâm $\widehat{BOC} = 120^\circ$. Tiếp tuyến của đường tròn $(O)$ tại $B$ và $C$ cắt nhau tại điểm $A$. Đoạn thẳng $OA$ cắt dây cung $BC$ tại điểm $H$.
\begin{enumerate}[a)]
	\item Chứng minh $OA$ là đường trung trực của đoạn thẳng $BC$ và bốn điểm $A, B, O, C$ cùng thuộc một đường tròn.
	\item Kẻ đường kính $CE$ của đường tròn $(O)$, tia $AE$ cắt $(O)$ tại điểm thứ hai $D$ ($D \ne E$). Chứng minh $AC^2 = AD \cdot AE$.
	\item Tính diện tích hình viên phân giới hạn bởi dây cung $BC$ và cung nhỏ $BC$ của đường tròn $(O)$ (lấy $\pi \approx 3{,}14$, kết quả làm tròn đến chữ số thập phân thứ hai).
\end{enumerate}
\loigiai{
\begin{enumerate}[a)]
	\item Theo tính chất của hai tiếp tuyến cắt nhau, ta có $AB = AC$ và $OB = OC = R$, do đó $OA$ là đường trung trực của đoạn thẳng $BC \implies OA \perp BC$ tại $H$ và $H$ là trung điểm của $BC$.\\
	Vì $AB, AC$ là hai tiếp tuyến nên $\widehat{ABO} = 90^\circ$ và $\widehat{ACO} = 90^\circ$.\\
	Tứ giác $ABOC$ có $\widehat{ABO} + \widehat{ACO} = 90^\circ + 90^\circ = 180^\circ$, suy ra bốn điểm $A, B, O, C$ cùng thuộc đường tròn đường kính $OA$.
	\item Tam giác $CDE$ nội tiếp đường tròn $(O)$ có đường kính $CE$ nên $\widehat{CDE} = 90^\circ \implies CD \perp AE$ tại $D$.\\
	Xét tam giác $ADC$ và tam giác $ACE$ có:
	$$ \widehat{CAD} \text{ là góc chung}, \quad \widehat{ACD} = \widehat{AEC} \text{ (góc tạo bởi tiếp tuyến và dây cung cùng chắn cung } CD). $$
	Suy ra $\Delta ADC \backsim \Delta ACE$ (g-g) $\implies \dfrac{AC}{AE} = \dfrac{AD}{AC} \implies AC^2 = AD \cdot AE$.
	\item Diện tích hình quạt tròn $OBC$ chắn cung nhỏ $BC$ có số đo $120^\circ$ là:
	$$ S_{\text{quạt}} = \dfrac{\pi R^2 \cdot 120^\circ}{360^\circ} = \dfrac{3{,}14 \cdot 3^2}{3} = 9{,}42\text{ (cm}^2\text{)}. $$
	Diện tích tam giác $OBC$ là:
	$$ S_{\Delta OBC} = \dfrac{1}{2} \cdot OB \cdot OC \cdot \sin \widehat{BOC} = \dfrac{1}{2} \cdot 3 \cdot 3 \cdot \sin 120^\circ = \dfrac{9\sqrt{3}}{4} \approx 3{,}90\text{ (cm}^2\text{)}. $$
	Diện tích hình viên phân giới hạn bởi dây $BC$ và cung nhỏ $BC$ là:
	$$ S = S_{\text{quạt}} - S_{\Delta OBC} \approx 9{,}42 - 3{,}90 = 5{,}52\text{ (cm}^2\text{)}. $$
\end{enumerate}
}
\end{bt}
